\chapter{Computational Models of Motor Control}
\label{ch:computational}

\epigraph{%
  A theory of motor control must ultimately be a theory of
  computation: what is computed, how, and why.
}{}

\section{Optimal Feedback Control}

The dominant computational theory of motor control is optimal feedback control (OFC)
\cite{TodorovJordan2002}:

\begin{definition}{The OFC Framework}{ofc-def}
  The nervous system controls movement by solving, in real-time, the stochastic
  optimal control problem:
  \begin{equation}
    \min_{\control(\cdot)} \mathbb{E}\left[
    \phi(\state_T) + \int_0^T \ell(\state_t, \control_t) \dt
    \right],
    \label{eq:ch10:ofc}
  \end{equation}
  subject to:
  \begin{equation}
    \dd\state = f(\state, \control) \dt + G(\state) \dd\bm{w},
  \end{equation}
  where $\dd\bm{w}$ is Brownian noise representing signal-dependent motor noise.
\end{definition}

\noindent
This is precisely the optimal control problem studied in Agrachev \& Sachkov \cite{AgrachevSachkov2004},
Chapters~5--6 (stochastic optimal control). Their framework establishes existence and uniqueness
of optimal feedback laws, characterizes them via the Hamilton-Jacobi-Bellman equation, and shows
that feedback synthesis via optimal control naturally leads to closed-loop stability.

\subsection{The Minimum Intervention Principle}

OFC naturally explains why the nervous system tolerates task-irrelevant variability
(the UCM from Chapter~1): the optimal controller only corrects deviations that
affect the cost function. Task-irrelevant deviations are ``free'' and need not be
corrected.

\begin{lstlisting}[language=Python, caption={Linear-quadratic optimal feedback controller.}]
import numpy as np
from scipy.linalg import solve_continuous_are

def lqr_ofc(
    A: np.ndarray,
    B: np.ndarray,
    Q: np.ndarray,
    R: np.ndarray,
    state_noise_cov: np.ndarray | None = None
) -> tuple[np.ndarray, np.ndarray]:
    """Compute LQR optimal feedback gains.

    This implements the infinite-horizon OFC solution.

    Parameters
    ----------
    A : np.ndarray, shape (n, n)
        State dynamics matrix.
    B : np.ndarray, shape (n, m)
        Input matrix.
    Q : np.ndarray, shape (n, n)
        State cost matrix.
    R : np.ndarray, shape (m, m)
        Control cost matrix.
    state_noise_cov : np.ndarray, shape (n, n), optional
        Process noise covariance.

    Returns
    -------
    tuple
        (K, P) where K is the optimal gain and P is the Riccati solution.
    """
    # Solve continuous algebraic Riccati equation
    P = solve_continuous_are(A, B, Q, R)

    # Optimal gain
    K = np.linalg.solve(R, B.T @ P)

    return K, P


# Example: 2D reaching with motor noise
# State: [x, y, dx, dy], Control: [Fx, Fy]
A = np.array([
    [0, 0, 1, 0],
    [0, 0, 0, 1],
    [0, 0, -2, 0],  # damping
    [0, 0, 0, -2],
])
B = np.array([
    [0, 0],
    [0, 0],
    [1, 0],
    [0, 1],
])

# Task: reach to target with minimum effort
# High cost on position at target, low cost on path
Q = np.diag([10, 10, 1, 1])  # Position matters more than velocity
R = 0.01 * np.eye(2)  # Low control cost

K, P = lqr_ofc(A, B, Q, R)
print("OFC gain matrix K:")
print(K)
print("\nMinimum intervention: gains are higher for")
print("task-relevant (position) than task-irrelevant dimensions")
\end{lstlisting}

\section{Signal-Dependent Noise}

A key feature of biological motor systems is that motor noise is
\textbf{signal-dependent}: the noise variance scales with the control signal magnitude:
\begin{equation}
  \control = \bar{\control} + \sigma(\bar{\control}) \bm{\xi},
  \quad \sigma(\bar{\control}) = k \norm{\bar{\control}},
  \label{eq:ch10:signal-dep-noise}
\end{equation}
where $\bm{\xi}$ is white noise and $k$ is the noise coefficient.

This explains minimum-jerk trajectories: the optimal trajectory under signal-dependent
noise minimizes the time-integral of squared jerk \cite{FlashHogan1985, Harris1998}.

\section{Active Inference (Friston Framework)}

Active inference reframes motor control as inference: the agent infers the
actions that would make its predictions come true:
\begin{equation}
  F = \mathbb{E}_q[\ln q(\state) - \ln p(\state, \bm{y})],
  \label{eq:ch10:free-energy}
\end{equation}
where $F$ is the variational free energy, $q(\state)$ is the approximate posterior,
and $p(\state, \bm{y})$ is the generative model.

\section*{Exercises}

\begin{enumerate}
  \item \textbf{OFC Reaching.} Using the LQR code, simulate reaching movements
        to different targets. Show that task-irrelevant variability is tolerated.

  \item \textbf{Signal-Dependent Noise.} Simulate reaching with and without
        signal-dependent noise. Plot the trajectory variability as a function of
        movement speed.

  \item \textbf{(Programming.)} Implement the minimum-jerk model and the
        minimum-variance model. Show that both predict smooth bell-shaped velocity
        profiles.
\end{enumerate}
